\section{Relation to established results}
\label{sec:related-work}
The earlier QED models draw on established, regime-specific methods:
constant-field electric-probe response with a stated finite subtraction
\cite{mages2010euler}, homogeneous mean-field pair production with
adiabatic current renormalization \cite{kluger1993pair}, and locally
covariant current and stress renormalization with explicit background
restrictions \cite{zahn2014dirac}. These sources do not themselves
provide the workbench's evaluated finite-mode formulas or a general
strong-field Bianchi-I closure.

The lattice and Hamiltonian frameworks originate with Wilson and
Kogut--Susskind \cite{wilson1974confinement,kogut1975hamiltonian}.
For the particular reductions studied here, the physical single-plaquette
problem already has a Mathieu-function formulation
\cite[Appendix B]{bauer2023basis}, and the open two-plaquette electric
spectrum has explicit low-level tabulations
\cite[Appendix C]{froland2025two}.
Maximal-tree gauge fixing also has a developed operator treatment: its
electric generators require transported derivatives, beyond the kinematic
identification of gauge orbits \cite[Appendix C]{burbano2024gauge}.
Thus these reductions and low levels are comparison targets, not claims
of discovery. The possible increment is the explicit graph arithmetic,
Gram identities, complete residual accounting, and inspectable bounds
recorded below. The Gegenbauer polynomials and variational tools used in
those calculations are classical
\cite{nistDLMFGegenbauer,teschl2014methods}.

Product-vacuum stability already permits infinite-dimensional onsite
spaces under appropriate gap and bounded-interaction hypotheses
\cite{yarotsky2004quasiparticles}. Local stability of ground-state
expectations is a related statement whose local perturbation need not
preserve the perturbed global gap \cite[Theorem 3]{henheik2022local}.
Explicit qubit thresholds \cite{bravyi2008polynomial} and finite-dimensional
local block diagonalization \cite{delvecchio2022local} have narrower
hypotheses than some models here. However, a broader construction does
cover separable onsite spaces and relatively form-bounded interactions
\cite{delvecchio2021unbounded}. Its relevance prevents any claim that
unbounded-interaction stability is absent from the literature. Applying
it quantitatively still requires matching domains, interaction geometry,
weighted norms, and smallness constants. This compendium does not supply
an evaluated general homogeneous threshold by citing that theorem.

Locality and approximation likewise have established frameworks.
Nachtergaele--Sims treat unbounded onsite terms with controlled bounded
intersite interactions \cite{nachtergaele2014dynamics}; their operator
topologies must be retained when comparing time continuity and spatial
limits. Exponential-in-volume tails offer a further lead under a
finite-qudit setup \cite{mcdonough2025volume}, not an admitted extension
to these rotor factors. Rigorous quantum-number truncation estimates
\cite{tong2022simulation} require their specified initial-state and
Hamiltonian assumptions. They do not automatically certify an arbitrary
retained block or its numerically centered heat evolution.

Finally, strong-coupling Wilson-measure and stochastic-generator results
\cite{shen2022strong} concern explicitly specified measures and clocks.
They do not identify an auxiliary Markov gap with the canonical
physical-time gap. Continuum reconstruction requires a full compatible
Schwinger-function hierarchy with the corrected regularity and other
axioms \cite{osterwalder1975axioms}; the four-dimensional existence and
physical mass-gap target remains distinct \cite{jaffe_witten}.
This is a selected literature comparison, not an exhaustive priority
search. Scientific novelty of the model-specific certificates has not
been established.
